\akapit  Głównym celem ćwiczenia laboratoryjnego było zapoznanie się z podstawowymi własnościami systemów sterowania bezzałogowego statku powietrznego (BSL) oraz ocena wpływu niedoskonałości elementów toru sterowania i pomiaru na osiągi i właściwości pilotażowe modelu. W ramach prac badawczych przeanalizowano naturalną dynamikę obiektu w locie swobodnym, a także zbadano wpływ dynamiki serwomechanizmów, luzów kinematycznych, opóźnień transportowych oraz szumów pomiarowych na dokładność realizacji zadanych trajektorii lotu.
